\documentclass[11pt,a4paper]{article}
\usepackage[margin=2.2cm]{geometry}
\usepackage{amsmath,amssymb,xcolor,enumitem,booktabs}
\usepackage[normalem]{ulem}
\usepackage[hidelinks]{hyperref}
\setlength{\parskip}{0.6em}\setlength{\parindent}{0pt}
\definecolor{newc}{HTML}{1f6f43}
\newcommand{\new}[1]{\textcolor{newc}{#1}}
\newcommand{\old}[1]{\textcolor{gray}{\sout{#1}}}
\newenvironment{greenblock}{\color{newc}}{}
\title{Proposed change, the two legs of a shelf's trip}
\date{23 November 2026, version 3 (follows the code, equations fit the page)}
\begin{document}\sloppy
\maketitle

\section*{How to read this}
Version 2 follows the decisions of 23 November. The text follows the code,
because every run was trained and tested on it. The idea is said in plain
words first, then ``Formally,'' and the formula. No new name such as
``weighted trip'' is added. The two legs get their own symbols, because the
formula cannot be written right without them. Fixed storage's overflow gets
one sentence. The SKU with no order line is left out, because it never
happens in the runs.

Black text stays, \new{green} is new and \old{struck} text goes. Part~1 is
what the code does. Parts~2 to~4 show the paragraphs that change. Part~5 is
what changes elsewhere. Nothing below is in the thesis yet.

\section*{Part 1. What the code does}
Checked in \texttt{slap-mapd-coupling} on \texttt{origin/main},
\texttt{storage/arrival.py} (\texttt{\_refill\_scores},
\texttt{leg\_weighted\_distance}, \texttt{putaway\_fixed},
\texttt{putaway\_congestion}), with the run configs (pallets of 10,
put-away by the agents, docks drawn at random, categories).
\begin{itemize}[nosep]
  \item The dock leg is the mean distance from the docks to the shelf. The
        workstation leg is the mean distance from the shelf to the four
        workstations of the SKU's category. The text has both right.
  \item The legs are not added with the same weight. The workstation leg is
        weighted by the order lines of the SKU per step, one pick per line.
        The dock leg is weighted by the pallets of the SKU per step, the units
        ordered per step divided by the pallet size, one put-away per pallet.
  \item Congestion-aware storage adds $\beta$ times the standing still of the
        shelf to the workstation leg, so it is weighted by the lines per step
        too. This matches the blocking cost $K$.
  \item Demand and congestion-aware storage take the SKUs from most to least
        ordered. Fixed storage takes them in index order, home shelf first.
        A pallet that misses its home shelf goes by the same two weighted
        legs. For one SKU the order of its shelves depends only on how the
        two weights compare, about pallet size over units per line, $10/3$,
        the same for every SKU. So fixed storage favours no SKU.
  \item Equal scores are broken by the plain trip of Appendix~D, the mean
        over the docks plus the nearest workstation, then by vertex number.
  \item How often a pallet misses its home shelf under fixed storage, at 16
        agents: about 8 percent of the arriving units under Token Passing, 40
        to 60 percent under the learned controllers (long runs, test seeds),
        whose fleets fall behind and leave the shelves full.
\end{itemize}

\section*{Part 2. \texttt{4.Implementation.tex}, Where arriving stock goes}

\subsection*{The trip paragraph}
A shelf is judged by the trip a unit makes through it, from a dock to the
shelf and on to a workstation of the SKU's category. The first leg is the
mean over the docks, because the dock of a pallet is drawn at random. The
second leg is the mean over the category's four workstations, because an
order may use any of them.
\old{Both legs count timesteps, so they add without a weight. This is the
\emph{trip distance} $d^{\mathrm{trip}}(k,v)$ (C.2) (Appendix C.1).}
\begin{greenblock}
The two legs are not walked equally often. Every pick of the SKU goes from
the shelf to a workstation of its category, so the second leg counts once per
pick. Every pallet goes from a dock to the shelf, so the first leg counts once
per pallet. A line asks for 3 units on average and a pallet holds 10, so a
SKU is picked about three times as often as it gets a pallet, and the second
leg weighs more.

Formally, write $d^{\mathrm{dock}}(v)$ for the first leg and
$d^{\mathrm{work}}(k,v)$ for the second, both in timesteps (Appendix~C).
Write $\hat\rho^{\mathrm{pal}}_t(k)$ for the pallets of $k$ per step, the
units of $k$ ordered per step in the last storage period divided by the
pallet size $Q$ (``pal'' for pallet). Each leg is weighted by how often it is
walked,
\begin{equation*}
  \hat\rho_t(k)\, d^{\mathrm{work}}(k,v) + \hat\rho^{\mathrm{pal}}_t(k)\, d^{\mathrm{dock}}(v) .
\end{equation*}
Take tea with $\hat\rho_t = 0.3$ lines per step, so about $0.09$ pallets per
step. Shelf $A$ is 4 steps from tea's workstations and 10 from the docks.
Shelf $B$ is 8 and 4. Unweighted, $B$ is shorter, 12 against 14. Weighted,
$A$ scores $0.3\cdot4+0.09\cdot10=2.1$ and $B$ scores
$0.3\cdot8+0.09\cdot4=2.76$, so tea goes to $A$, closer to its workstations,
because it is picked far more often than it is put away.
\end{greenblock}
The demand rules are not told a SKU's category $h(k)$. They learn it from
its first order line, which always goes to a workstation of that category.
Fixed storage uses $h(k)$ directly.

\subsection*{The cost paragraph and eq.~(storageobjective)}
The demand rules can be read as keeping a cost small. The classic slotting
cost $D$ adds up, over the pallets placed at an update, how often their SKU
is \old{ordered times the length of their trip}\new{picked times the second
leg, plus how often it gets a pallet times the first leg}. The blocking cost
$K$ is the thesis's own [\dots unchanged]. The two costs and the value the
rules aim to keep small are
\begin{align*}
D &= \sum_{k\in\mathcal{K}}\sum_{v\in V_{\mathrm{str}}}
  \bigl(\hat\rho_t(k)\,d^{\mathrm{work}}(k,v)+\hat\rho^{\mathrm{pal}}_t(k)\,d^{\mathrm{dock}}(v)\bigr)\,x^{\mathrm{new}}_t(k,v),\\
K &= \sum_{k\in\mathcal{K}}\sum_{v\in V_{\mathrm{str}}}
  \hat\rho_t(k)\,\hat\omega_t(v)\,x^{\mathrm{new}}_t(k,v),
  \qquad D+\beta K .
\end{align*}
The old first line was
$D=\sum_k\sum_v \old{\hat\rho_t(k)\,d^{\mathrm{trip}}(k,v)}\,x^{\mathrm{new}}_t(k,v)$.
\new{In words, the bracket is the picks times the second leg, plus the
pallets times the first leg.}

\subsection*{Fixed storage}
[\dots] SKUs are taken in index order. A pallet that does not fit on the home
shelf goes to the nearest shelf \old{by $d^{\mathrm{trip}}(k,v)$} with room
for it\new{. ``Nearest'' counts both legs, the trip to the workstations
about three times as much as the trip from the docks, as for every rule. That
weight comes from how often a SKU is picked against how often it gets a
pallet, which is about the same for every SKU. So the overflow shelf does not
depend on how popular the SKU is, and fixed storage still favours no SKU}.

\subsection*{Demand storage}
[\dots] It takes the SKUs from most to least ordered, by the estimated demand
$\hat\rho_t(k)$, and gives each the \old{nearest shelves by
$d^{\mathrm{trip}}(k,v)$}\new{shelves with the smallest weighted legs. The
weights differ per SKU, so each SKU has its own order of shelves}. This keeps
the slotting cost $D$ small, with $\beta=0$, greedily and not exactly.

\subsection*{Congestion-aware storage, eq.~(storagegreedy)}
[\dots] Formally, the \emph{congestion-aware score} of SKU $k$ at shelf $v$ is
\begin{equation*}
  \mathrm{score}(k,v) = \hat\rho_t(k)\bigl(d^{\mathrm{work}}(k,v) + \beta\,\hat\omega_t(v)\bigr)
  + \hat\rho^{\mathrm{pal}}_t(k)\,d^{\mathrm{dock}}(v) .
\end{equation*}
The old score was $\hat\rho_t(k)\bigl(d^{\mathrm{trip}}(k,v)+\beta\,\hat\omega_t(v)\bigr)$.
\new{Standing still joins the second leg, because every pick from the shelf
meets it.} [\dots rest unchanged]

\subsection*{The ties sentence}
Ties, between SKUs or between shelves, are broken by index, SKU first and
then shelf, \new{and equal scores first by the plain trip of
Appendix~D,} so each rule gives exactly one placement.

\section*{Part 3. Appendix C, the trip distance (\texttt{eq:tripdistance})}
\old{The trip distance,}
\begin{equation*}
d^{\mathrm{trip}}(k,v) = d^{\mathrm{dock}}(v) + d^{\mathrm{work}}(k,v) \quad\text{(old, removed)}
\end{equation*}
\new{The two legs of a shelf's trip,}
\begin{greenblock}
\begin{equation*}
  d^{\mathrm{dock}}(v) = \frac{1}{|V_{\mathrm{dock}}|} \sum_{g \in V_{\mathrm{dock}}} d_G(g, v),
  \qquad
  d^{\mathrm{work}}(k,v) = \frac{1}{|h(k)|} \sum_{w \in h(k)} d_G(v, w),
\end{equation*}
\end{greenblock}
where $h(k)$ is the category of SKU $k$ [\dots unchanged].
\new{A pallet that misses its home shelf under fixed storage is placed by
these legs. At 16 agents this is about 8 percent of the arriving units under
Token Passing and 40 to 60 percent under the learned controllers, whose
fleets fall behind and leave the shelves full.}

\section*{Part 4. Problem Formulation}
No change. It says demand storage picks ``the shelves with the shortest
trip'', which stays true in words.

\section*{Part 5. Elsewhere in the same change}
\begin{itemize}[nosep]
  \item \texttt{A.Notation.tex}: add $d^{\mathrm{dock}}(v)$,
        $d^{\mathrm{work}}(k,v)$ and $\hat\rho^{\mathrm{pal}}_t(k)$. No trip
        distance row exists today.
  \item The label \texttt{eq:tripdistance} stays on the Appendix~C equation,
        so references still work. Its name becomes ``the two legs''.
  \item Appendix~D keeps its plain trip $d^{\mathrm{trip}}(v)$. It describes
        the warehouse and breaks ties, as the code uses it.
  \item The worked example of tea is in the text. It can also be drawn on
        the tea figure of the Problem Formulation (\texttt{fig:storageexample})
        once the text is approved.
  \item \texttt{GAPS.md}: a new entry, text follows the code, no model
        change. Tags on the changed lines.
\end{itemize}

\end{document}
